\section{ELA-CNN: Training Details and Additional Results}\label{app:cnn}

\textbf{Selection and training.} The network was trained on the train split (8,828 images) with horizontal flips,
class-weighted binary cross-entropy, AdamW (learning rate $3\cdot10^{-4}$, weight decay $10^{-4}$), batch 64 and a
cosine schedule over at most 15 epochs; the epoch and pooling rule were selected by AUC on the validation split (1,894
images). Early stopping ended training after 4 epochs without improvement in validation AUC;
nothing besides the epoch and the pooling rule was tuned, because a second 5-fold CV, as used for the classical models,
would have cost about $15\times20$~min of GPU time. Table~\ref{tab:app-cnn-sel} lists the selected recipes and
Fig.~\ref{fig:app-cnn-training} the validation curves. Training and refitting took about 3.5~hours on an Apple M4 GPU
(MPS backend), which is not bit-reproducible. A preflight step checked the models, the ELA settings, the classical
scores and every cached comparison result before the single test run, whose start and completion are recorded in the
run log. Code: \texttt{scripts/make\_ela.py}, \texttt{scripts/run\_phase8.py}, \texttt{scripts/phase8\_posthoc.py} and
\texttt{src/forgery/cnn.py}.

\textbf{ELA input.} Fig.~\ref{fig:cnn-ela} shows the network's input for two validation images, as released and after
protocol R. The authentic JPEG shows almost no ELA error as released, while the tampered TIFF shows strong error along
edges and texture: the format leak is visible in ELA itself.

\begin{figure}[!ht]
\centering
\includegraphics[width=0.7\textwidth]{fig_ela_examples.png}
\caption{ELA inputs of the CNN for two validation images. Top: an authentic JPEG, whose ELA as released (A) is almost
black. Bottom: a tampered TIFF, whose ELA as released shows strong error along edges and texture. Right: ELA of the
same images after protocol R.}
\label{fig:cnn-ela}
\end{figure}


\begin{table}[!ht]
\centering
\caption{Selected epoch and pooling rule per protocol (train split; validation split for selection).}
\label{tab:app-cnn-sel}
\footnotesize
\begin{tabular}{lccl}
\toprule
Protocol & Selected epoch & Pooling & Validation AUC \\
\midrule
R & 3 & max & 0.746 (seeds 42 / 43 / 44: 0.746 / 0.762 / 0.752) \\
B & 8 & mean & 0.797 \\
A & 12 & mean & 0.994 \\
\bottomrule
\end{tabular}
\end{table}

\begin{figure}[!ht]
\centering
\includegraphics[width=0.6\textwidth]{fig_cnn_training.png}
\caption{Validation AUC (best pooling) per epoch of the ELA-CNN under A, B and R; under R for three seeds.}
\label{fig:app-cnn-training}
\end{figure}

\textbf{Image size.} Table~\ref{tab:app-cnn-size} extends the post-hoc size grouping of Section~\ref{sec:cnn-groups}.
The train-only selection models show the same drop on validation: R 0.728 on standard-size vs.\ 0.567 on larger
images; B 0.792 vs.\ 0.664. Under B, larger authentic images also receive higher scores than standard-size authentic
ones (mean logit 0.17 vs.\ $-$1.76). Within the larger images, the score correlates only weakly with the number of
tiles (Spearman 0.20 under R, max pooling; $-$0.09 under B, mean pooling). Against the forest on forensic + shortcut
features, the CNN differs by $-$0.045\ci{$-$0.074}{$-$0.017} (R), $-$0.125\ci{$-$0.145}{$-$0.104} (B) and
$-$0.005\ci{$-$0.009}{$-$0.003} (A). Under B, CNN $-$ classical is $-$0.129\ci{$-$0.152}{$-$0.106} on copy-move and
$-$0.071\ci{$-$0.101}{$-$0.042} on splicing.

\begin{table}[!ht]
\centering
\caption{Test AUC by natural image size (post-hoc; the pre-registered area terciles are invalid).}
\label{tab:app-cnn-size}
\footnotesize
\begin{tabular}{lcccccc}
\toprule
Size (test images) & $n$ & Tampered & R: CNN & R: classical & R: CNN $-$ classical & B: CNN $-$ classical \\
\midrule
Smaller than $384\times256$ & 21 & 62\% & 0.702 & 0.615 & +0.087\ci{$-$0.072}{0.260} & $-$0.144 \\
$384\times256$, either orientation & 1,469 & 34\% & 0.770 & 0.789 & $-$0.019\ci{$-$0.065}{0.026} & $-$0.105\ci{$-$0.133}{$-$0.080} \\
Larger than $384\times256$ & 402 & 65\% & 0.544 & 0.730 & $-$0.186\ci{$-$0.232}{$-$0.139} & $-$0.217\ci{$-$0.266}{$-$0.163} \\
\bottomrule
\end{tabular}
\end{table}

\textbf{Robustness.} Table~\ref{tab:app-cnn-robust} gives the test AUC after each degradation of
Section~\ref{sec:robustness}; each was followed by the protocol and then ELA (classical results are the frozen outputs
of that section). Under R, the CNN loses about as much as the classical pipeline: JPEG Q75 costs 0.008 vs.\ 0.006,
resize $\times$0.5 0.085 vs.\ 0.112, and the social upload 0.037 vs.\ 0.044; after resizing, the gap between the two
is not significant ($-$0.008\ci{$-$0.043}{0.027}). Under B, the classical pipeline loses much more, because its B
advantage rests on compression traces that any re-save overwrites: JPEG Q75 costs it 0.180 against 0.065 for the CNN,
leaving the two level (+0.007\ci{$-$0.034}{0.047}), and the social upload 0.199 against 0.137. Padding is not the cause of
the CNN's losses: after downscaling most images become smaller than one tile (95\% under R resize $\times$0.5; 82\%
after a social upload), and scoring them without padding changes the AUC by at most 0.018. On MICC-F220 the CNN reaches
0.529\ci{0.360}{0.700} (R) and 0.454\ci{0.226}{0.660} (B), consistent with chance, while the classical detector
reaches 0.972 (R) and 0.902 (B); the CIs are wide because the forgeries come from only 11 scenes. The failure fits the
size analysis: all MICC-F220 photographs are larger than CASIA's standard size ($722\times480$ to $800\times1{,}070$)
and come from a different dataset, whereas the classical copy-move matching finds the duplicated regions regardless of
their source.

\begin{table}[!ht]
\centering
\caption{Test AUC of the ELA-CNN and the classical pipeline under degradations and on MICC-F220.}
\label{tab:app-cnn-robust}
\footnotesize
\begin{tabular}{lccccc}
\toprule
Condition & R: CNN & R: classical & R: rank average & B: CNN & B: classical \\
\midrule
clean & 0.770 & 0.805 & 0.824 & 0.811 & 0.920 \\
JPEG Q75 & 0.762 & 0.799 & 0.817 & 0.746 & 0.739 \\
resize $\times$0.5 & 0.685 & 0.693 & 0.720 & 0.672 & 0.698 \\
social upload & 0.733 & 0.762 & 0.788 & 0.674 & 0.721 \\
MICC-F220 & 0.529\ci{0.360}{0.700} & 0.972 & 0.811 & 0.454\ci{0.226}{0.660} & 0.902 \\
\bottomrule
\end{tabular}
\end{table}
